% 第4章 4.1.1：三种测试函数列，psi(x)=exp(-1/(1-x^2))。
% (a) psi/n（n=1,2,4）在测试意义下趋零；(b) phi_n=psi sin(12x)/12 的值（虚线，振幅<=0.031）很小而导数（实线，最大约0.37）不小；
% (c) psi(x-3k)/k（k=1,2,3）值与导数都趋零，但支撑跑向无穷。数据由同目录 generate.py 计算。
\begin{tikzpicture}[font=\figurefont\footnotesize]
\pgfplotsset{every axis/.style={width=50mm,height=34mm,axis lines=left,
  x axis line style={draw=FigureStroke,line width=.6pt},y axis line style={draw=FigureStroke,line width=.6pt},
  tick style={draw=FigureStroke},tick label style={font=\figurefont\scriptsize,text=FigureInk},
  yticklabel style={/pgf/number format/fixed},
  xlabel={$x$},xlabel style={at={(axis description cs:1,0)},anchor=west,font=\figurefont\scriptsize},
  title style={font=\figurefont\footnotesize,align=center},
  legend style={draw=none,fill=none,font=\figurefont\scriptsize,cells={anchor=west},at={(0.5,-0.3)},anchor=north,legend columns=3}}}
\def\dat{figures/02-abstract-mathematics/03-function-spaces/tikz/c4-test-sequences}
\begin{axis}[name=A,xmin=-1.3,xmax=1.3,ymin=0,ymax=0.42,xtick={-1,0,1},ytick={0,0.2,0.4},title={(a) 支撑固定，导数趋零}]
\addplot[draw=FigureBlue,line width=1.1pt] table[x=x,y=n1] {\dat/a.dat};
\addplot[draw=FigureGreen,line width=1.1pt,dashed] table[x=x,y=n2] {\dat/a.dat};
\addplot[draw=FigureYellow,line width=1.1pt,dash dot] table[x=x,y=n4] {\dat/a.dat};
\legend{$n=1$,$n=2$,$n=4$}
\end{axis}
\begin{axis}[name=B,at={(A.east)},anchor=west,xshift=13mm,xmin=-1.3,xmax=1.3,ymin=-0.42,ymax=0.42,xtick={-1,0,1},ytick={-0.4,0,0.4},title={(b) 值趋零，导数不趋零}]
\addplot[draw=FigureBlue,line width=1.1pt] table[x=x,y=dphi] {\dat/b.dat};
\addplot[draw=FigureYellow,line width=1.1pt,dashed] table[x=x,y=phi] {\dat/b.dat};
\legend{$\varphi_{12}'$,$\varphi_{12}$}
\end{axis}
\begin{axis}[at={(B.east)},anchor=west,xshift=13mm,xmin=-0.5,xmax=10.5,ymin=0,ymax=0.42,xtick={0,3,6,9},ytick={0,0.2,0.4},title={(c) 值趋零，支撑逃逸}]
\addplot[draw=FigureBlue,line width=1.1pt] table[x=x,y=k1] {\dat/c.dat};
\addplot[draw=FigureGreen,line width=1.1pt,dashed] table[x=x,y=k2] {\dat/c.dat};
\addplot[draw=FigureYellow,line width=1.1pt,dash dot] table[x=x,y=k3] {\dat/c.dat};
\legend{$k=1$,$k=2$,$k=3$}
\end{axis}
\end{tikzpicture}
